% ============================================================
%  CHAPTER 5 — Implementation
% ============================================================

\chapter{Implementation}
\label{ch:implementation}

\section{System Overview}

The firmware maps the software architecture described in
Chapter~\ref{ch:architecture} directly onto the ESP32's dual-core execution model.
At boot, \texttt{setup()} initialises LittleFS, loads NVS configuration, initialises the sensors, starts
both Wi-Fi interfaces, registers HTTP routes, and launches the Firebase FreeRTOS task .
After \texttt{setup()} returns, the Arduino \texttt{loop()} continues to execute, calling \texttt{updateSensors()} on every iteration and scheduling periodic
Wi-Fi background scans.

% -------------------------------------------------------
% TODO (Backend Team): describe how the backend connects to
% and consumes from Firebase (or your chosen database),
% how the server is started, and how it links to the
% architecture described in Chapter 3.
% -------------------------------------------------------
\begin{center}
	\textit{[Backend system overview to be added.]}
\end{center}

% -------------------------------------------------------
% TODO (Web Dashboard Team): describe how the web application
% connects to the backend and how it relates to the
% component architecture described in Chapter 3.
% -------------------------------------------------------
\begin{center}
	\textit{[Web dashboard system overview to be added.]}
\end{center}

% ============================================================
\section{Hardware Implementation}
% ============================================================

The hardware assembly consists of three components mounted on a common 3.3\,V power
rail:

\begin{itemize}
	\item \textbf{ESP32 development board}.
	\item \textbf{DHT22} connected to GPIO~4.
	\item \textbf{MAX30102 breakout board} connected to the ESP32's default
	      I\textsuperscript{2}C bus (SDA: GPIO~21, SCL: GPIO~22), powered at 3.3\,V.
\end{itemize}

All components are powered from the ESP32 development board's 3.3\,V regulated output.
The board is powered from USB 5\,V connected to a mains adapter in the hospital room.

% ============================================================
\section{Hardware Schematics}
% ============================================================

Figure~\ref{fig:schematic} provides the complete wiring diagram of the assembled circuit.


\begin{figure}[H]
	\centering
	\includegraphics[width=\linewidth]{hardware_diagram.png}
	\caption{Hardware wiring diagram}
	\label{fig:schematic}
\end{figure}

% ============================================================
\section{Software Development Environment}
% ============================================================

Table~\ref{tab:dev_env} lists the complete development environment used for the
ESP32 embedded subsystem.

\begin{table}[H]
	\centering
	\caption{ESP32 software development environment.}
	\label{tab:dev_env}
	\begin{tabularx}{\linewidth}{lX}
		\toprule
		\textbf{Component} & \textbf{Details}                               \\
		\midrule
		Operating System   & Linux (Debian 13 trixie)                       \\
		ESP32 Arduino core v3.3.7                                           \\
		Board target       & \texttt{esp32:esp32:esp32da} (ESP32-WROOM-32)  \\
		Flash tool         & \texttt{esptool} (Python); \texttt{mklittlefs} \\
		Serial monitor     & Custom \texttt{monitor.sh}                     \\
		\bottomrule
	\end{tabularx}
\end{table}

% -------------------------------------------------------
% TODO (Backend Team): describe your development environment
% (language, framework, runtime version, database engine,
% deployment platform, IDE, dependency manager).
% -------------------------------------------------------
\begin{center}
	\textit{[Backend development environment to be added.]}
\end{center}

% -------------------------------------------------------
% TODO (Web Dashboard Team): describe your development
% environment (Node version, framework, build tool, etc.).
% -------------------------------------------------------
\begin{center}
	\textit{[Web dashboard development environment to be added.]}
\end{center}

% ============================================================
\section{Software Implementation}
% ============================================================

\subsection{Embedded Subsystem}

\subsubsection{Firebase Initialisation on Core~0}

The \texttt{initFirebase()} function is called from inside
\texttt{firebaseUploadTask()}, which is pinned to Core~0 via
\texttt{xTaskCreatePinnedToCore(..., 0)}.

\subsubsection{Configuration Portal}

The configuration portal is served entirely from the LittleFS partition.
The \texttt{isAPConnected()} guard checks whether the incoming request arrived via
the AP subnet (\texttt{WiFi.softAPIP()}) and redirects any STA-side request to the
root path, preventing access from the hospital network.

\subsubsection{Wi-Fi Recovery}

All reconnection logic is consolidated in the \texttt{wifiConnect()} helper, which
always calls \texttt{WiFi.disconnect()} followed by
\texttt{WiFi.begin(ssid, password)}.

% -------------------------------------------------------
% TODO (Backend Team): describe your backend implementation
% details (route handlers, middleware, database queries,
% alert engine implementation, authentication logic).
% -------------------------------------------------------
\subsection{Backend Implementation}
\begin{center}
	\textit{[To be completed by the backend team.]}
\end{center}

% -------------------------------------------------------
% TODO (Web Dashboard Team): describe the web application
% implementation (key components, real-time update logic,
% chart implementation, alert display logic).
% -------------------------------------------------------
\subsection{Web Dashboard Implementation}
\begin{center}
	\textit{[To be completed by the web dashboard team.]}
\end{center}

% ============================================================
\section{Software Validation}
% ============================================================

\subsection{Embedded Subsystem Validation}

Validation was performed across three dimensions:

\textbf{Sensor output validation.}
DHT22 readings were cross-referenced against a calibrated reference thermometer.

\textbf{Firebase upload validation.}
The Firebase Realtime Database console was monitored in real time to confirm record
arrival, correct path structure (\texttt{/rooms/\{roomId\}} and
\texttt{/patients/\{patientId\}}), correct field names and data types, and
correct server-generated timestamps.

\textbf{Authentication validation.}
The serial monitor (debug build) was used to trace the complete authentication event
sequence: \textit{authenticating} $\to$ \textit{auth request sent} $\to$
\textit{auth response received} $\to$ \textit{token obtained}.
Error conditions (invalid API key, incorrect password) were deliberately induced
and confirmed to produce the expected HTTP~400 error codes via the callback.

% -------------------------------------------------------
% TODO (Backend Team): describe your API validation
% protocol (unit tests, integration tests, Postman
% collections, or equivalent).
% -------------------------------------------------------
\subsection{Backend Validation}
\begin{center}
	\textit{[To be completed by the backend team: API test protocol, test cases,
				and results.]}
\end{center}

% -------------------------------------------------------
% TODO (Web Dashboard Team): describe UI validation
% (browser tests, manual test scenarios, responsive
% design validation on target devices).
% -------------------------------------------------------
\subsection{Web Dashboard Validation}
\begin{center}
	\textit{[To be completed by the web dashboard team.]}
\end{center}

% ============================================================
\section{Embedded Implementation (ESP32)}
% ============================================================

\subsection{Memory Optimisation}

The LittleFS image partition is mounted with \texttt{LittleFS.begin(true)}, enabling automatic
formatting on first boot.
NVS storage uses the \texttt{Preferences} library with the \texttt{"settings"}
namespace, consuming at most 10 string entries at typical credential lengths.

\subsection{Latency Optimisation}

Delta filtering (Equations~\ref{eq:room_thresh} and~\ref{eq:patient_thresh}) reduces
the frequency of HTTPS requests to only those carrying new information, keeping the
per-upload round-trip time as the dominant latency term rather than queuing delays.

% -------------------------------------------------------
% TODO (AI Team): if deploying any model on the ESP32 or
% on the cloud, describe the model/algorithm adaptation,
% memory optimisation, latency optimisation, and any
% modifications made to fit the platform.
% -------------------------------------------------------
\section{AI Integration}
\begin{center}
	\textit{[To be completed by the AI team: model conversion, deployment mechanism
				(cloud API / TFLite / ONNX), integration with the data pipeline, and known
				limitations.]}
\end{center}
